% ============================================================================
% Section 5.4 — Solving Multi-Step Problems with Rational Numbers
% CCSS 7.EE.B.3
% ============================================================================
\section{Solving Multi-Step Problems with Rational Numbers}

\begin{stepGoal}
\begin{itemize}
    \item Solve equations with fraction and decimal coefficients
    \item Use estimation to check that answers are reasonable
\end{itemize}
\end{stepGoal}

\begin{stepVocabBox}
    \vocabItem{LCD}{Least Common Denominator --- the smallest number all denominators divide into.}
    \vocabItem{Clearing Fractions}{Multiplying every term by the LCD to remove fractions from an equation.}
\end{stepVocabBox}

\begin{stepsCard}[How to Solve Equations with Fractions or Decimals]
    \stepItem{\textbf{Estimate} the answer first so you can check reasonableness later.}
    \stepItem{If the equation has fractions, \textbf{multiply every term by the LCD} to clear them. If it has decimals, work with them directly.}
    \stepItem{\textbf{Solve} the simpler equation using two-step methods.}
    \stepItem{\textbf{Check}: substitute back and compare to your estimate.}
\end{stepsCard}

% --- Worked Example 1 (clearing fractions) ---
\begin{stepExample}{Solve $\dfrac{2x}{5} - 3 = 7$.}
    \stepShow{1} Estimate: $\frac{2x}{5}$ must equal $10$, so $x$ is around $25$.

    \stepShow{2} Add $3$: $\dfrac{2x}{5} = 10$. Multiply both sides by $5$: $2x = 50$.

    \stepShow{3} Divide by $2$: $x = 25$.

    \stepShow{4} Check: $\frac{2(25)}{5} - 3 = 10 - 3 = 7$ \checkmark \quad Matches estimate.
    \stepResult{$x = 25$}
\end{stepExample}

% --- Worked Example 2 (decimals) ---
\begin{stepExample}{Solve $0.3x + 2.1 = 5.1$.}
    \stepShow{1} Estimate: $0.3x$ needs to equal about $3$, so $x \approx 10$.

    \stepShow{2} Decimals are fine --- subtract $2.1$: $0.3x = 3.0$.

    \stepShow{3} Divide by $0.3$: $x = 10$.

    \stepShow{4} Check: $0.3(10) + 2.1 = 3 + 2.1 = 5.1$ \checkmark
    \stepResult{$x = 10$}
\end{stepExample}

\watchOut{When clearing fractions, multiply \textbf{every} term by the LCD --- not just the fraction term. Missing a term is the most common mistake!}

% ============================================================================
% PRACTICE
% ============================================================================
\newpage
\begin{stepPractice}{Multi-Step Problems with Rational Numbers Practice}
    \resetProblems

    \stepPracticeHeader[step1Color]{Estimate First}
    \prob For $0.5x + 3 = 10$, estimate $x$ before solving. \answerBlank[2cm]
    \answer{$x \approx 14$}

    \stepPracticeHeader[step2Color]{Clear Fractions or Work with Decimals}
    \begin{multicols}{2}
    \prob Solve $\dfrac{x}{4} + \dfrac{3}{4} = 5$. \answerBlank[2cm]
    \answer{$x = 17$}
    \prob Solve $-1.5y + 4 = -8$. \answerBlank[2cm]
    \answer{$y = 8$}
    \end{multicols}

    \stepPracticeHeader[step3Color]{Solve and Check}
    \prob Solve $\dfrac{3n}{2} - 5 = 10$. \answerBlank[2cm]
    \answer{$n = 10$}

    \prob Solve $0.25x + 1.5 = 4$. \answerBlank[2cm]
    \answer{$x = 10$}

    \wordProblem{A store marks down shirts by $\frac{1}{3}$ of the price, then takes off another $\$4$. The final price is $\$12$. What was the original price?}{dollars}
    \answer{$\$24$}
\end{stepPractice}
